\chapter{کلاس NP}

برای مسئله $\prob{PATH}$ توانستیم با دوری از جست‌وجوی کامل، الگوریتمی چندجمله‌ای پیدا کنیم. اما برای بعضی مسئله‌های دیگر، که بسیاری از آن‌ها جالب و کاربردی هم هستند، تا امروز هیچ الگوریتم چندجمله‌ای‌ای پیدا نشده است.

\section{یک مثال: مسیر همیلتونی}

\begin{definition}[\lr{HAMPATH}]
یک \textbf{مسیر همیلتونی} در گراف جهت‌دار $G$ مسیری جهت‌دار است که از هر رأس \textbf{دقیقاً یک بار} عبور کند.
\[\prob{HAMPATH}=\{\enc{G,s,t}\st G\ \text{\rl{گرافی جهت‌دار با یک مسیر همیلتونی از $s$ به $t$ است}}\}.\]
\end{definition}
\begin{ltrfig}
\begin{tikzpicture}[->,>={Stealth[round]},semithick,every node/.style={circle,draw=jzMain,fill=jzMain!7,minimum size=5mm,inner sep=0pt,font=\small}]
  \node (s) at (0,0) {$s$}; \node (a) at (1.4,1) {}; \node (b) at (1.4,-1) {}; \node (c) at (2.8,0.35) {};
  \node (d) at (2.8,-0.35) {}; \node (e) at (4.2,1) {}; \node (f) at (4.2,-1) {}; \node (t) at (5.6,0) {$t$};
  \draw (s)--(a); \draw (s)--(b); \draw (b)--(a); \draw (a)--(c); \draw (a)--(e); \draw (c)--(d);
  \draw (c)--(e); \draw (d)--(b); \draw (d)--(f); \draw (f)--(e); \draw (e)--(t); \draw (f)--(t);
\end{tikzpicture}
\end{ltrfig}
در این گراف مسیر همیلتونی‌ای از $s$ به $t$ وجود دارد (آن را پیدا کنید!).

ما نمی‌دانیم چرا نتوانسته‌ایم برای بعضی مسئله‌ها الگوریتم چندجمله‌ای پیدا کنیم. شاید چنین الگوریتم‌هایی وجود دارند و هنوز پیدا نشده‌اند؛ شاید هم این مسئله‌ها ذاتاً سخت‌اند. اما یک نکته مهم را می‌دانیم: پیچیدگی بسیاری از این مسئله‌ها به هم \textbf{وابسته} است؛ یک الگوریتم چندجمله‌ای برای یکی از آن‌ها می‌تواند کل یک کلاس از مسئله‌ها را حل کند.

این مسئله‌ها یک ویژگی مشترک دارند که \textbf{وارسی‌پذیری چندجمله‌ای} نام دارد: پیدا کردن جواب شاید سخت باشد، اما \textbf{بررسی} یک جواب پیشنهادی آسان است. مثلاً اگر کسی یک مسیر همیلتونی از $s$ به $t$ به ما بدهد، به‌راحتی (در زمان چندجمله‌ای) می‌توانیم بررسی کنیم که درست است.

\section{وارسی‌کننده و کلاس NP}

\begin{definition}[وارسی‌کننده]
یک \term{وارسی‌کننده}{Verifier} برای زبان $A$ الگوریتمی مثل $V$ است که
\[A=\{w\st V\ \text{\rl{رشته $\enc{w,c}$ را برای یک رشته $c$ می‌پذیرد}}\}.\]
زمان وارسی‌کننده را فقط بر حسب طول $w$ می‌سنجیم؛ پس یک \textbf{وارسی‌کننده با زمان چندجمله‌ای} در زمان چندجمله‌ای بر حسب طول $w$ اجرا می‌شود. زبان $A$ \textbf{وارسی‌پذیر چندجمله‌ای} است اگر یک وارسی‌کننده با زمان چندجمله‌ای داشته باشد.
\end{definition}

\begin{itemize}
\item به $c$ یک \textbf{گواه} (\lr{certificate}) یا \textbf{اثبات عضویت} در $A$ می‌گوییم.
\item چون وارسی‌کننده در زمان چندجمله‌ای اجرا می‌شود، فقط می‌تواند بخشی چندجمله‌ای از گواه را بخواند؛ پس می‌توان فرض کرد طول گواه چندجمله‌ای (بر حسب طول $w$) است.
\item برای $\prob{HAMPATH}$ گواه همان مسیر همیلتونی از $s$ به $t$ است.
\end{itemize}

\begin{note}
بعضی مسئله‌ها ممکن است وارسی‌پذیر چندجمله‌ای نباشند. مثلاً متمم مسئله، $\overline{\prob{HAMPATH}}$: حتی اگر بدانیم گراف مسیر همیلتونی \emph{ندارد}، نمی‌دانیم چطور بدون یک الگوریتم نمایی این را به کسی ثابت کنیم؛ باید \emph{نبودن} چیزی را وارسی کرد.
\end{note}

\begin{definition}[کلاس NP]
$\cNP$ کلاس زبان‌هایی است که \textbf{وارسی‌کننده با زمان چندجمله‌ای} دارند.
\end{definition}

نام $\cNP$ از \textbf{زمان چندجمله‌ای غیرقطعی} (\lr{Nondeterministic Polynomial time}) آمده است. دلیلش این است که روش دیگری برای توصیف همین کلاس وجود دارد: مسئله‌هایی که با یک ماشین تورینگ \textbf{غیرقطعی} در زمان چندجمله‌ای حل می‌شوند.

\begin{example}[\lr{HAMPATH} با ماشین غیرقطعی]
\begin{algo}{$N_1$}{روی ورودی $\enc{G,s,t}$ که $G$ گرافی جهت‌دار با رأس‌های $s$ و $t$ است}
\begin{enumerate}
\item فهرستی از $m$ عدد $p_1,\dots,p_m$ بنویس که $m$ تعداد رأس‌های $G$ است. هر عدد به‌طور غیرقطعی بین $1$ و $m$ انتخاب می‌شود.
\item اگر عدد تکراری در فهرست بود، \textbf{رد کن}.
\item اگر $s\neq p_1$ یا $t\neq p_m$، \textbf{رد کن}.
\item برای هر $i$ بین $1$ و $m-1$ بررسی کن که $(p_i,p_{i+1})$ یالی از $G$ است. اگر نبود \textbf{رد کن}؛ اگر همه بودند \textbf{بپذیر}.
\end{enumerate}
\end{algo}
هر مرحله در زمان چندجمله‌ای اجرا می‌شود (مرحله ۱ در هر شاخه $O(m\log m)$ و بقیه هم چندجمله‌ای)، پس $N_1$ یک ماشین غیرقطعی با زمان چندجمله‌ای است.
\end{example}

\begin{theorem}\label{thm:np-ntm}
یک زبان در $\cNP$ است اگر و تنها اگر یک ماشین تورینگ غیرقطعی با زمان چندجمله‌ای آن را تصمیم بگیرد.
\end{theorem}

\begin{proofidea}
یک وارسی‌کننده با زمان چندجمله‌ای را به یک ماشین غیرقطعی با زمان چندجمله‌ای تبدیل می‌کنیم و برعکس. ماشین غیرقطعی وارسی‌کننده را با \textbf{حدس زدن گواه} شبیه‌سازی می‌کند. وارسی‌کننده هم ماشین غیرقطعی را با استفاده از \textbf{شاخه پذیرنده به عنوان گواه} شبیه‌سازی می‌کند.
\end{proofidea}

\begin{proof}
\textbf{از وارسی‌کننده به ماشین غیرقطعی:} فرض کنید $A\in\cNP$ و $V$ یک وارسی‌کننده با زمان $n^k$ برای $A$ باشد. ماشین $N$:
\begin{algo}{$N$}{روی ورودی $w$ به طول $n$}
\begin{enumerate}
\item به‌طور غیرقطعی رشته‌ای مثل $c$ با طول حداکثر $n^k$ انتخاب کن.
\item $V$ را روی ورودی $\enc{w,c}$ اجرا کن.
\item اگر $V$ پذیرفت \textbf{بپذیر}؛ وگرنه \textbf{رد کن}.
\end{enumerate}
\end{algo}
$V$ در زمان $n^k$ حداکثر $n^k$ نماد از گواه را می‌خواند، پس محدود کردن طول $c$ به $n^k$ چیزی را از دست نمی‌دهد. $N$ رشته $w$ را می‌پذیرد اگر و تنها اگر گواهی مثل $c$ باشد که $V$ بپذیرد، یعنی اگر و تنها اگر $w\in A$. هر شاخه $N$ در زمان چندجمله‌ای تمام می‌شود.

\textbf{از ماشین غیرقطعی به وارسی‌کننده:} فرض کنید ماشین غیرقطعی $N$ در زمان چندجمله‌ای $A$ را تصمیم بگیرد. وارسی‌کننده $V$:
\begin{algo}{$V$}{روی ورودی $\enc{w,c}$ که $w$ و $c$ رشته هستند}
\begin{enumerate}
\item $N$ را روی ورودی $w$ شبیه‌سازی کن و هر نماد $c$ را توصیف انتخاب غیرقطعی در یک گام در نظر بگیر (مثل آدرس‌دهی گره‌ها در اثبات قضیه \ref{thm:ntm}).
\item اگر این شاخه از محاسبه $N$ پذیرفت \textbf{بپذیر}؛ وگرنه \textbf{رد کن}.
\end{enumerate}
\end{algo}
اگر $w\in A$، شاخه پذیرنده‌ای وجود دارد و وقتی گواه درست (انتخاب‌های آن شاخه) داده شود، $V$ دقیقاً آن شاخه را اجرا می‌کند و می‌پذیرد. اگر $w\notin A$، هیچ شاخه‌ای نمی‌پذیرد، پس $V$ برای هیچ $c$ای نمی‌پذیرد. $V$ فقط یک شاخه را اجرا می‌کند، پس در زمان چندجمله‌ای است.
\end{proof}

\begin{definition}[کلاس زمانی غیرقطعی]
\[\mathrm{NTIME}(t(n))=\{L\st L\ \text{\rl{زبانی است که یک ماشین تورینگ غیرقطعی با زمان $O(t(n))$ تصمیمش می‌گیرد}}\}.\]
\end{definition}

\begin{corollary}
$\cNP=\bigcup_k\mathrm{NTIME}(n^k)$.
\end{corollary}

\section{پرسش P در برابر NP}

به زبان ساده:
\begin{itemize}
\item $\cP$ = کلاس زبان‌هایی که عضویت در آن‌ها را می‌توان به‌سرعت \textbf{تصمیم گرفت}.
\item $\cNP$ = کلاس زبان‌هایی که عضویت در آن‌ها را می‌توان به‌سرعت \textbf{وارسی کرد}.
\end{itemize}

روشن است که $\cP\subseteq\cNP$: اگر بتوانیم جواب را سریع پیدا کنیم، بدون نیاز به گواه هم می‌توانیم آن را سریع بررسی کنیم. به نظر می‌رسد وارسی‌پذیری چندجمله‌ای خیلی قوی‌تر از تصمیم‌پذیری چندجمله‌ای باشد. اما هیچ‌کس تا امروز نتوانسته \textbf{ثابت کند} مسئله‌ای وجود دارد که وارسی‌کننده چندجمله‌ای دارد ولی تصمیم‌گیرنده چندجمله‌ای ندارد.

\begin{ltrfig}
\begin{tikzpicture}[font=\small]
  \draw[thick,draw=jzMain,fill=jzMain!6,rotate=20] (0,0) ellipse (1.2 and 1.7); \node at (0.45,1.0) {NP};
  \draw[thick,draw=jzAccent,fill=jzAccent!12] (-0.3,-0.55) circle (0.65); \node at (-0.3,-0.55) {P};
  \draw[thick,draw=jzMain,fill=jzMain!6] (4,0) circle (1.0); \node at (4,0) {P = NP};
  \node at (2.05,0) {or};
\end{tikzpicture}
\end{ltrfig}

پرسش «آیا $\cP=\cNP$؟» یکی از بزرگ‌ترین مسئله‌های حل‌نشده علوم کامپیوتر نظری است. بیشتر پژوهشگران باور دارند $\cP\neq\cNP$. این مسئله یکی از هفت «مسئله هزاره» مؤسسه ریاضی کلی است که برای حل هر کدام یک میلیون دلار جایزه گذاشته شده است:
\begin{itemize}
\item حدس برچ و سوینرتون-دایر،
\item حدس هاج،
\item وجود و همواری جواب معادلات ناویر–استوکس،
\item \textbf{مسئله P در برابر NP}،
\item فرضیه ریمان،
\item وجود نظریه یانگ–میلز و شکاف جرم،
\item حدس پوانکاره (که تنها مسئله حل‌شده این فهرست است).
\end{itemize}

\begin{summary}
\begin{itemize}
\item وارسی‌کننده $V$: $w\in A$ اگر و تنها اگر گواهی $c$ باشد که $V$ رشته $\enc{w,c}$ را بپذیرد.
\item $\cNP$: زبان‌های دارای وارسی‌کننده چندجمله‌ای $=$ زبان‌هایی که ماشین غیرقطعی در زمان چندجمله‌ای تصمیمشان می‌گیرد $=\bigcup_k\mathrm{NTIME}(n^k)$.
\item $\prob{HAMPATH}\in\cNP$ (گواه: خود مسیر). $\cP\subseteq\cNP$ و پرسش $\cP=\cNP$ هنوز باز است.
\end{itemize}
\end{summary}
